%--------------------------------------------------------------------%
% Page          : Appendix A - Experiment Configuration
%--------------------------------------------------------------------%

\chapter*{LAMPIRAN A: Konfigurasi Eksperimen}
\addcontentsline{toc}{chapter}{LAMPIRAN A: Konfigurasi Eksperimen}

\section*{Parameter Kalibrasi}

Tabel berikut menyajikan parameter kalibrasi yang digunakan dalam eksperimen, bersumber dari \texttt{CalibrationConfig} (\texttt{libs/shared/shared/models/calibration.py}).

\begin{table}[htbp]
    \centering
    \begin{tabular}{lr}
        \toprule
        Parameter & Nilai \\
        \midrule
        fib\_n & 33 \\
        io\_wait\_ms & 50,0 \\
        lambda\_compute\_ms & 24,0 \\
        r\_saturation\_per\_replica & 66,7 \\
        target\_cpu\_util & 0,5 \\
        min\_k8s\_replicas & 3 \\
        max\_k8s\_replicas & 10 \\
        pod\_cpu\_millicores & 300 \\
        node\_cpu\_limit & 1,0 \\
        proactive\_trend\_threshold & 3,0 \\
        proactive\_approach\_ratio & 0,6 \\
        proactive\_lookahead\_steps & 5 \\
        \bottomrule
    \end{tabular}
\end{table}

\section*{Alokasi Sumber Daya Kluster}

\begin{table}[htbp]
    \centering
    \begin{tabular}{lrr}
        \toprule
        Komponen & CPU (Docker) & Pod/Replica \\
        \midrule
        K8s server node & 3,0 & --- \\
        K8s agent node & 3,0 & --- \\
        Serverless agent node & 3,0 & --- \\
 Dynamic workload nodes & 1,0 each & --- \\
        K8s baseline pods & --- & 3 $\times$ 300m \\
        Knative pods (KPA) & --- & 3 (min-scale) \\
        \bottomrule
    \end{tabular}
\end{table}

\section*{Parameter \emph{Trace} ClarkNet}

\begin{table}[htbp]
    \centering
    \begin{tabular}{lr}
        \toprule
        Parameter & Nilai \\
        \midrule
        Jumlah tahap & 40 \\
        Durasi per tahap & 30 detik \\
        Total durasi & 1.200 detik (20 menit) \\
        RPS minimum & 22 \\
        RPS maksimum & 164 \\
        RPS rata-rata & 73 \\
        Jumlah \emph{request} per \emph{run} & 87.840 \\
        \bottomrule
    \end{tabular}
\end{table}
